\chapter{The Bracketed Number}

\section{The bound is the object}

The slip arrived as a bracket, and the machine learned to read in brackets. This chapter takes the bracket seriously as
the \emph{object} the instrument owns --- not a stand-in for a point it is still hunting, but the thing itself, complete.
The machine owns the interval. It does not own, and does not claim, the point inside it. This section is about taking
that claim literally, because the whole ending depends on the machine meaning it.

The ordinary picture has the bracket as scaffolding: a value is really a point, and the interval is a temporary
confession of ignorance the machine will eventually collapse. This machine rejects that picture. For a value it cannot
compute exactly, there \emph{is} no point the machine possesses --- there is only the interval it cornered, and the
interval is the whole of what it has. Interval arithmetic makes this a discipline rather than an apology: a quantity is
carried as a pair, a lower bound and an upper, and every operation on it produces another pair, so that the machine
never at any stage holds a point it did not earn. The bound is not a box around the answer. The bound is the answer, in
the only form a finite machine can honestly hold one.

This inverts what precision means, and the inversion is worth dwelling on. Naively, a wider bracket is a worse answer
and the goal is to shrink it to nothing. But a bracket shrunk below what the machine can certify is not a better answer;
it is a lie, a point asserted where only an interval was earned. The honest goal is not the narrowest bracket but the
\emph{true} one --- the interval exactly as wide as the machine's certified computation makes it, neither padded nor
pinched. A machine that owns its interval reports that true width and stops. It does not chase the point, because the
point is not among the things it has; chasing it would mean fabricating certainty. Owning the bound means being content
with the interval and refusing the point --- which is precisely the contentment the last chapter of the book will need.

There is a name in this book for what sits inside the bound and is not the machine's to name: the point is the
\emph{generic} value, approached by the bracket and captured by no finite condition, exactly the shape the number
chapter gave a real. The machine's interval is a finite condition; the point it brackets is the generic limit no finite
condition reaches. So owning the bound and disowning the point is not modesty --- it is the plain structure of what a
number is to a finite machine. The bracket is the finite thing the machine holds; the point is the generic thing it
approaches; and the honest instrument reports the first and refuses to counterfeit the second.

In the two verbs the bound is a funge the machine declines to close. It funges together, into one interval, every value
its two computed walls cannot decide apart --- pools the indistinguishable into a single bag. And it leaves the bag
open: it does not tange out, from among the pooled values, a single point to call the answer, because no computation it
ran decided one. The bound is the open funge owned as an object. Where a lesser instrument would pull the bag shut on a
point, this one holds it open, exactly as wide as its decided walls, and calls the open bag the number. What sets the
width of the bag --- the floor below which the walls cannot be told apart --- is the next section's.

\section{The floor}

The bracket is the machine's object, and its width is set by one thing: the \emph{floor}. Below a certain resolution the
machine cannot tell two candidates apart, and where distinction fails, bracketing must stop. This section builds that
floor exactly, because it is the single most important quantity in the whole instrument --- the depth at which the
machine's honesty and its limitations become the same thing. The floor is where the machine runs out of decidable
distinctions, and everything the book refuses to claim, it refuses because of the floor.

Every finite representation has one, and numerical analysis has a precise name for it: \emph{machine epsilon}, the
smallest gap the representation can resolve, the distance below which two values round to the same
mark~\cite{golub2013}. Add to a number something smaller than machine epsilon and the number does not change --- the
addition falls beneath the floor and registers as nothing. This is not a defect of a particular representation; it is a
property of \emph{any} finite one. A machine with finitely many symbols can name only finitely many values, so between
any two adjacent nameable values there is a gap, and inside that gap the machine is blind. The floor is that blindness,
made exact: the resolution below which the machine's distinctions all return \emph{indistinguishable}.

Now the bisection meets the floor, and the meeting sets the bracket's width. The machine halves its interval, tightening
the walls, each halving a decidable comparison. But the comparisons are themselves computed to the machine's
resolution, and once the walls fall within machine epsilon of each other, the comparison that would tighten them further
returns \emph{same} --- the two walls have become one value to the machine, and there is nothing left to decide. The
bisection halts there, not because the value is pinned but because the machine has exhausted the distinctions it can
make. The bracket's width is therefore not arbitrary and not a choice; it is the floor, surfaced --- the interval is as
wide as the smallest distinction the machine can still certify, and no narrower.

And here is the honesty in the floor, which is the whole of why the book insists on it. Below the floor there is still
\emph{something} --- a sub-floor tail, the part of the value finer than any distinction the machine can draw. That tail
does not vanish; the machine simply cannot see it. The honest thing is to say so: the sub-floor tail is the machine's
own ignorance, named and bounded, not a region the machine claims is empty. A dishonest machine would keep halving past
the floor, manufacturing distinctions its resolution cannot support, and report a point --- pretending the sub-floor
tail was zero when it only could not be seen. The honest machine halts at the floor and hands back the bracket, with the
sub-floor tail acknowledged as unresolved. The floor is where the machine's power ends, and marking it exactly is how
the machine tells the truth about where that is.

In the two verbs the floor is where tange runs out. Above it, the machine tanges --- decides candidates apart, tightens
the walls, splits the interval finer. At the floor the tange fails: the comparison returns \emph{indistinguishable}, and
the machine can no longer decide one candidate from another. So it funges --- pools all the sub-floor candidates into
the bracket and stops, because below the floor there is no difference left to tange. The floor is the exact depth at
which the machine's first verb gives out and only its second remains: it cannot tell the sub-floor values apart, so it
bags them together and holds the bag open. That open bag, floored, is the honest number --- and the last thing this
chapter must show is why the machine can be trusted to have drawn its floor honestly, which is repeatability.

\section{Repeatability is the honesty}

The machine owns a bracket floored at its resolution. But an interval on a page is only as trustworthy as the machine's
ability to produce it \emph{again}. This closing section names the property that makes a bracket a reading rather than an
assertion: \emph{repeatability}. A reading is not merely a bracket the machine computed once. It is a bracket the machine
can run again and get back the same. Repeatability is the honesty --- the thing that separates a measurement from a
lucky number.

Start with what repeatability demands. Run the machine's bracketing procedure once, note the walls; run it again from
the same start, and the walls must come back identical --- same lower, same upper, to the last resolvable mark. This is
determinism, and for a finite machine it is a checkable property: the computation has no hidden state, no unrepeatable
step, nothing that would let a second run differ from the first. A bracket that came back different on a rerun would not
be a reading; it would be noise wearing a reading's clothes, and the machine could not stand behind it. The machine
stands behind its brackets precisely because they are reproducible: what it computed, it can compute again, and the
agreement of the runs is the evidence that the bracket is real.

This is the computational form of the oldest rule in measurement, that a result you cannot reproduce is not yet a
result. The machine keeps the rule in its strictest form: a reading is defined as \emph{the bracket that reruns to
itself}. If two runs disagree, the disagreement is above the floor and the machine has not yet reached its resolution ---
it must keep bracketing until the runs agree, and where they agree is the floor. So repeatability and the floor are the
same fact seen twice: the floor is the depth at which further distinctions fail, and repeatability is the guarantee that
at that depth the machine returns the same bracket every time. A reading is where the machine's runs converge and stay
converged --- the interval that no rerun can budge.

And this is why the ending can be trusted, planted here on small readings so it is ready for the large one. The coupling
at the end will be a bracket, floored at the machine's resolution, and the reason to believe it is not that the machine
asserts it but that the machine can produce it again --- and, beyond that, that four independent gauges each produce their
own version of it and those versions agree. Repeatability within a gauge is the first cross-check; agreement across
gauges is the second. Both are the same principle: a reading is what survives being taken again. The machine earns the
right to hand back a bracket instead of a point by proving the bracket reproducible, and it learns to prove that here,
on the floor of its own resolution, long before the number that matters.

In the two verbs repeatability is a funge across runs. The machine tanges each run's bracket apart --- computes it as
its own object --- and then funges the runs together: it pools the first run's bracket with the second's, and when they
decide \emph{same}, the pooled thing is the reading, one interval that two runs agree on. A bracket that could not be
funged with its own rerun --- that came back different --- would be no reading at all. So the honest number is a funge
of a value with itself across repetition: the interval that reruns to the same interval, floored, owned, and
reproducible. With the bracket now a genuine reading, the machine is ready to ask what such a reading \emph{is} --- and
the next part discovers that one bracket has more than one face.
